\documentclass[hyphens,aspectratio=169,dvipsnames]{beamer}
\usepackage{graphicx}
\usepackage{xcolor}
\usepackage{amssymb}
\usepackage{amsmath}
\usepackage{mathtools}
\usepackage{textcomp}
\usepackage{moresize}
\usepackage{framed}
\usepackage{minted}
\usepackage{relsize}
\usepackage{tikz}
\usepackage{comment}
\usetikzlibrary{shapes.geometric, arrows, positioning}
\usetheme{Berlin}
\usecolortheme[RGB={0,95,47}]{structure}
\beamertemplatenavigationsymbolsempty

\makeatletter
\AtBeginEnvironment{minted}{\dontdofcolorbox}
\def\dontdofcolorbox{\renewcommand\fcolorbox[4][]{##4}}
\makeatother

\newcommand{\textpf}[1]{\texttt{\color{black}\fcolorbox{lightgray}{lightgray}{#1}}}

\begin{document}

\begin{frame}{Learning objectives}
    \begin{itemize}
        \item Sign extension
        \item Bitwise operations
    \end{itemize}
    
\end{frame}

\begin{frame}{2\'s complement}
        Recall how signed numbers are represented in binary: 2\'s complement
        \begin{itemize}
            \item To negate a number, flip all the bits and add one
        \end{itemize}
\end{frame}

\section{Sign extension}

\begin{frame}[fragile]{Example Program}
    What does this program do?
    \begin{center}
        \begin{minted}{c}
#include <stdint.h>
#include <stdio.h>

int main(void) {
    int8_t i8 = -1;
    int16_t i16 = i8;
}
        \end{minted}
    \end{center}
\end{frame}

\begin{frame}{Sign Extension}
    \begin{itemize}
        \pause \item When converting a negative integer from a smaller bitwidth to a larger bitwidth, you need to fill the new bits with \texttt{1}s in order to preserve the original integer's value.
    \end{itemize}

    \pause For example,
    \begin{align*}
        \texttt{(int16\_t)(int8\_t)(-1)} &= \texttt{(int16\_t)0b11111111} \\
                                         &= \texttt{0b}\underbrace{\texttt{11111111}}_\text{(8 new \texttt{1}s)}\texttt{11111111} \\
                                         &= \texttt{(int16\_t)(-1)}
    \end{align*}
\end{frame}

\begin{frame}{Sign Extension}
    \begin{itemize}
        \pause \item When converting a nonnegative integer from a smaller bitwidth to a larger bitwidth, you need to fill the new bits with \texttt{0}s in order to preserve the original integer's value.
    \end{itemize}

    \pause For example,
    \begin{align*}
        \texttt{(int16\_t)(int8\_t)(127)} &= \texttt{(int16\_t)0b01111111} \\
                                          &= \texttt{0b}\underbrace{\texttt{00000000}}_\text{(8 new \texttt{0}s)}\texttt{01111111} \\
                                          &= \texttt{(int16\_t)(127)}
    \end{align*}
\end{frame}


\end{document}
